\section{Investment (Ugo)}\label{sec:investment}

Investment carries the rule under test. It makes three decisions: how much capacity to add,
which technique the new plants embody, and how to pay for them. Only the second may differ
between the two types.

\paragraph{The trap: quantity and finance can settle the contest before the technique rule
does.} If firms invest all their cash, capacity grows with profits and the firms with the
higher profit rate at the going price, the Okishio firms, grow faster. Investment also
creates the excess capacity that starts price wars, which favour Shaikh firms. The volume
rule sets the balance between the two selection channels, so it is identical for both types.

\paragraph{Version 1 rule.}
\[
\qbar^{\,d}_{i,t} = (1+g)\, x_{i,t-1}, \qquad
I_{i,t} = \max\{0,\ \qbar^{\,d}_{i,t} - \qbar_{i,t}\}, \qquad
j^{*} = \arg\min_{j}\, \big[(1 + h_i\tau)\, ucs_j + (\delta + h_i)\, \kappa_j\big],
\]
with $g = 0$, $\tau = 0$, $\qbar_{i,t}$ the capacity left after scrapping plus plants on
order, and the menu C, D1, D2. Shaikh firms use $h = 0$, which picks the lowest unit cost;
Okishio firms use the best profit rate on the menu at their own current price, floored at
zero, which picks the highest profit rate at that price. New plants arrive next period and
are paid from cash, then debt up to twice the period's sales revenue; investment is scaled
down when it cannot be financed.

\paragraph{Version 2 rules (the defaults).} \code{PLANT\_SIZE} $= S$ gives every plant the
capacity $S$, including the initial plants (split into the nearest number of plants) and
entrants'. \code{INVESTMENT\_RULE = "retained\_earnings"} replaces the accelerator: whatever
the demand, firm $i$ buys as many plants of the technique $j^*$ as its retained earnings pay
for,
\[
I_{i,t} = S \left\lfloor \frac{\min\{RE_{i,t},\ M_{i,t}\}}{\kappa_{j^*} S} \right\rfloor,
\qquad RE_{i,t} = RE_{i,t-1} + \Pi_{i,t} - \mathit{Div}_{i,t} - \kappa_{j^*} I_{i,t-1},
\]
with $\Pi$ profit after depreciation and interest, $M$ cash, and $I$ all spending on plants,
replacement included. \code{BORROWING = False} removes borrowing for plants everywhere:
investment, under either rule, and replacement are paid from cash only, and the debt cap no
longer binds anything. With \code{CASH\_SHORTFALL = "bankruptcy"} there is no overdraft
either (section~\ref{sec:accounts}). Without a demand condition,
capacity is bounded only by scrapping: without \code{IDLE\_SCRAPPING} it grows without limit
(section~\ref{sec:scrapping}, \code{docs/design/choices.md}).

\paragraph{The two channels.} The retained-earnings rule is the accumulation channel of
the brief's trap, chosen deliberately: capacity grows with each type's profits, and the
firms with the higher profit rate at the going price grow faster. The price channel runs
alongside through market shares. The decomposition of each type's capacity share into net
investment, entry and exit, and of its sales share into price-driven, entry-exit and
rationing terms (\code{docs/design/collector-brief.md}, \S4.2--4.3;
\code{analysis/indicators.py}), is reported with every result. The version 1 accelerator
with borrowing (\code{INVESTMENT\_RULE = "accelerator"}, \code{BORROWING = True}) is the
other arm of the finance treatment.

\paragraph{Technical change (version 2 default).} With \code{TECH\_DRAWS = True}, the default (\code{TECH\_CONFLICT} $= q = 0.5$), a firm that invests is
offered one new design $n$, drawn relative to the technique $o$ it would otherwise build and
judged at its own price $p$. Every design is Marx-biased: operating cost falls by a share
$a \le$ \code{TECH\_STEP} and capital per unit rises by a share $b$,
\[
ucs_n = (1-a)\,ucs_o, \qquad \kappa_n = (1+b)\,\kappa_o .
\]
Unit cost falls if $b < b_S = a\,ucs_o/(\delta\kappa_o)$, and the profit rate at $p$ rises if
$b < b_O = a\,ucs_o/(\delta\kappa_o + p - uc_o)$. With probability \code{TECH\_CONFLICT} the
design lies in the conflict zone $b_O \le b < b_S$, which only the unit-cost rule adopts;
otherwise $0 < b < b_O$ and both rules adopt it. The conflict zone vanishes as the margin
$p - uc_o$ falls to zero, where the two rules coincide. The firm adopts $n$ if
$ucs_n + (\delta + h_i)\kappa_n < ucs_o + (\delta + h_i)\kappa_o$, the same comparison as the
technique choice, and keeps it among the techniques it knows; if not, it still invests, in
$o$, so a rejected design never costs a firm capacity. Each draw is logged (its zone, its
unit cost, whether it was adopted, and the order that builds it).

This rule has limits that keep it from carrying the contest on its own: the conflict share
$q$ is set by a coin flip rather than produced by an innovation set, so a price war cannot
close the conflict zone except at zero margin; no design is ever rejected by a unit-cost
firm; there are no capital-saving designs; the draw reads the firm's price; a firm that
invests more often is offered more designs; and designs are private to the adopter but
public to entrants. The innovation set below answers them; it is the second arm of the
switch \code{TECH\_CHANGE}, the default (\code{"categories"}) staying as it is until the
sensitivity analysis.

\paragraph{Innovation set} (\code{TECH\_CHANGE = "innovation\_set"},
\code{equations/innovation\_set.py}). A design $n$ is drawn relative to a reference
technique $o$ with nothing read from the firm but $o$:
\[
ucs_n = (1-a)\,ucs_o, \qquad \kappa_n = (1+b)\,\kappa_o, \qquad
a = u\,\text{\code{TECH\_STEP}}, \qquad b = k\,a, \qquad k = m\,e^{\sigma z},
\]
with $u$ uniform on $(0, 1]$ and $z$ standard normal (\code{TECH\_SET = "marx"}: every
design saves labour and uses capital). $k$ is the capital per unit given up per unit of
operating cost saved, log-normal around the median $m$ (\code{TECH\_BIAS} $= 1.5$, the
treatment) with spread $\sigma$ (\code{TECH\_BIAS\_SPREAD} $= 0.8$). The adoption
inequality is unchanged, so in terms of $k$ the unit-cost rule adopts if
$k < k_S = ucs_o/(\delta\kappa_o)$ and the profit-rate rule at price $p$ if
$k < k_O(p) = ucs_o/(\delta\kappa_o + \max\{0, p - uc_o\})$: the zones (both adopt,
conflict, both reject) are outcomes of the draw and of the price, and a price war closes
the conflict zone. For C at $p = 100$, $k_O = 1.80$ and $k_S = 4.67$; at the median bias of
the Penn World Table ($m = 0.53$) the rules agree on 94\,\% of draws at normal prices and
on all of them in a war; at $m = 1.5$ to $3$ the conflict share is 20 to 45\,\% at normal
prices and zero in a war (\code{docs/design/review-v2.md}, section 4, zone shares). The
neutral set (\code{TECH\_SET = "neutral"}) draws the direction uniformly on the circle in
the plane $(\Delta ucs/ucs, \Delta\kappa/\kappa)$ with size $u\,$\code{TECH\_STEP}:
capital-saving and dearer designs appear, the second wedge (only the profit-rate rule
adopts) with them, and realised Marx bias becomes an outcome, read from the
capacity-weighted $ucs$ and $\kappa$ of each type in the \code{groups} table. Three further
switches: the reference is the firm's best known technique (\code{TECH\_REFERENCE = "own"})
or the industry's lowest-unit-cost technique in use (\code{"frontier"}); designs are
offered once per firm and period with probability \code{TECH\_DRAW\_RATE} $= 0.2$ whether
or not the firm invests (\code{TECH\_DRAW\_TIMING = "period"}; an adopted design is built at
the next investment) or at each investment as in the categories rule
(\code{"investment"}); a design becomes known to every firm \code{IMITATION\_LAG} $= 2$
periods after its first plant produces (\code{TECH\_DIFFUSION = "public"}), or stays with
its adopter and later entrants (\code{"private"}). Every draw is logged with its zone at
the firm's price; the label swap holds under each switch.

\paragraph{Tests} (\code{tests/test\_investment.py}): at $p = 100$ Shaikh firms build D2 and
Okishio firms D1, and with only C and D2 on offer Okishio firms keep C; on random menus the
argmin with $h$ equal to the best profit rate picks the same technique as maximising the
profit rate directly; with constant demand and no scrapping, investment stops once capacity
meets expected sales; borrowing for plants never takes debt above its cap; two firms of
different types with the same sales add the same capacity. For the switches: whole plants
for a capacity; retained earnings buy whole plants and never borrow; with \code{BORROWING = False} no firm
borrows for plants over a run, under either investment rule, and replacement waits for cash; every plant, entrants'
included, has the fixed size; every design uses more capital and less operating cost and is
adopted exactly by the rules its category names; at zero margin both rules adopt a conflict
design; operating costs stay positive over a long run; a rejected design still leaves the
firm investing in its best known technique, an adopted one is built and kept; the label
swap holds with draws.

\subsection*{Code}
\eqfile{hurdle.py}
\eqfile{technique_choice.py}
\eqfile{desired_capacity.py}
\eqfile{investment.py}
\eqfile{retained_earnings_investment.py}
\eqfile{plant_number.py}
\eqfile{new_technique.py}
\eqfile{innovation_set.py}
\eqfile{borrowing_capacity.py}
\marketfile{plant_market.py}
\schedulerpart{investment}{\code{model/scheduler.py}: step 5}
